% ============================================================
%  chapters/ch13-gesture-reconstruction.tex
% ============================================================

\chapter{Gesture Reconstruction as Inverse Inference}

\chapterepigraph{%
  Recognition is not classification.\\
  Classification asks: which box?\\
  Recognition asks: what happened?
}{Flyxion, \textit{Semantic Infrastructure}}

\sidenote{Connects to\\Bayesian inverse\\problems and\\App.\,F §F.14}

Conventional gesture recognition systems are classifiers: given an input
(a signal, an image, a sequence of keystrokes), output a label from a
fixed vocabulary. This chapter argues that classification is the wrong
model. The right model is \emph{inverse inference}: given an observation
(a projection of the gesture), infer the most probable complete gesture
that produced it.

\section{Why Classification is Insufficient}

Classification assumes a fixed set of possible outputs and assigns
the input to the best-matching category. For gesture recognition,
this means: given a hand motion, output the most likely character, word,
or command from a predetermined vocabulary.

The problem with classification is that it is \emph{domain-closed}: it
can only recognize gestures in its training vocabulary. It cannot handle:

\begin{itemize}
  \item \textbf{Novel gestures:} gestures not in the training set,
        including new words, technical terms, proper nouns.
  \item \textbf{Degraded inputs:} gestures that are partially obscured,
        interrupted, or produced under unusual conditions.
  \item \textbf{Compositional meanings:} the meaning of a gesture
        sequence is not always the concatenation of individual gesture
        meanings (syntax matters).
  \item \textbf{Continuous signals:} the boundary between gesture classes
        is not sharp; classification forces a hard decision where
        the underlying signal is continuous.
\end{itemize}

Inverse inference handles all of these naturally, because it reasons
about the underlying gesture rather than the output category.

\section{Trajectory Reconstruction}

The core of inverse inference is trajectory reconstruction: given
observed data $O$ (a projection of the gesture), find the distribution
over complete gestures $\gamma$ that could have produced $O$.

\begin{definition}{Trajectory Reconstruction Problem}{traj-reconstruct}
  Given:
  \begin{itemize}
    \item Observation $O$ (projection of some unknown gesture),
    \item Forward model $P(O | \gamma)$ (likelihood of observing $O$
          given gesture $\gamma$),
    \item Prior $P(\gamma)$ (distribution over gestures in the
          admissible region).
  \end{itemize}
  The \emph{trajectory reconstruction} is the posterior distribution:
  \[
    P(\gamma | O) = \frac{P(O | \gamma) \cdot P(\gamma)}{P(O)}.
  \]
\end{definition}

This is a Bayesian inverse problem. The prior $P(\gamma)$ encodes
knowledge of the gesture manifold: admissible gestures have high
prior probability; inadmissible gestures have zero prior probability.
The likelihood $P(O | \gamma)$ encodes the projection model: how
likely is observation $O$ given that the underlying gesture was $\gamma$?

\begin{remark}
  In the Spherepop framework, the prior $P(\gamma)$ is determined by the
  accessibility entropy $S(\gamma)$: gestures in high-accessibility regions
  of the gesture manifold are more probable a priori, because they have
  more neighboring admissible gestures. The projection model
  $P(O | \gamma)$ is the probabilistic version of the canonical
  projection $\Phi \circ \hat{C}$ from Appendix~F.
\end{remark}

\section{Embodied Semantics}

The trajectory reconstruction framework provides a rigorous formalization
of \emph{embodied semantics}: the view that meaning is grounded in
bodily experience and motor behavior, not in abstract symbol manipulation.

\begin{definition}{Embodied Semantic Content}{embodied-sem}
  The \emph{embodied semantic content} of a gesture $\gamma$ is
  the posterior distribution $P(\gamma | O)$: the distribution over
  all possible gestures (including their full kinematic details)
  consistent with the observed projection $O$.
\end{definition}

The traditional symbol ``a'' has \emph{symbolic} semantic content:
it refers to a category of objects, sounds, or characters. The embodied
semantic content of the gesture that produced ``a'' includes all of this
plus the kinematic signature of the producer: their hand size, writing
style, emotional state, familiarity with the text, and motor history.

Embodied semantics is richer than symbolic semantics. It is the difference
between knowing that someone said ``hello'' and knowing \emph{how} they
said it — tentatively, warmly, anxiously, routinely. The word is the same;
the gesture is different; the embodied content is different.

\section{Constraint Multiplication}

Inverse inference becomes dramatically more powerful when multiple
observations provide simultaneous constraints on the unknown gesture.

\begin{definition}{Constraint Multiplication}{constraint-mult}
  Given $n$ independent observations $O_1, \ldots, O_n$, each
  constraining $\gamma$, the joint posterior is:
  \[
    P(\gamma | O_1, \ldots, O_n) \propto P(\gamma) \prod_{i=1}^n P(O_i | \gamma).
  \]
  Each additional observation \emph{multiplies} a constraint factor,
  sharpening the posterior.
\end{definition}

Constraint multiplication is the mathematical reason why multi-modal
sensing is powerful. A keyboard alone provides timing data but no spatial
data. A camera alone provides spatial data but misses key timing. Together,
they multiply constraints and dramatically narrow the posterior, making
reconstruction possible where either alone would be ambiguous.

\begin{example}
  Consider reconstructing the gesture for typing the word ``constraint.''
  \begin{itemize}
    \item Typing stream alone: provides key sequence and timing. The
          posterior over kinematic gestures is wide — many hand
          configurations could produce this timing.
    \item Camera alone: provides hand shape and position. The posterior
          over typed words is wide — many words could be produced with
          similar hand shapes.
    \item Together: timing constrains the temporal structure; camera
          constrains the spatial structure. Their product narrows
          the posterior sharply, and reconstruction becomes precise.
  \end{itemize}
\end{example}

\section{Latent Intention}

The deepest application of inverse inference is the reconstruction of
\emph{latent intention}: not just what gesture was made, but what the
person was \emph{trying} to do.

Intention is a higher-level description of the gesture — the goal that
drove the motor behavior, which may be only partially achieved. A person
trying to type ``receive'' who accidentally types ``recieve'' made a
gesture consistent with their intention but not their intention's result.

\begin{definition}{Latent Intention Model}{latent-intention}
  The \emph{latent intention} is a hidden variable $I$ in a hierarchical
  Bayesian model:
  \[
    P(O, \gamma, I) = P(O | \gamma) \cdot P(\gamma | I) \cdot P(I).
  \]
  Given observation $O$, the reconstructed intention is the marginal
  posterior $P(I | O) = \int P(I | \gamma) P(\gamma | O)\,d\gamma$.
\end{definition}

Autocorrect systems that correct spelling errors are primitive
implementations of latent intention reconstruction: they infer the
intended word from the actually typed word. More sophisticated systems
would infer the intended meaning from the intended word, allowing
correction of semantically incorrect but syntactically valid text.

\section{The Inverse Inference Principle}

\begin{namedthm}{The Inverse Inference Principle}
  Gesture recognition is not categorization from a fixed vocabulary.
  It is probabilistic reconstruction of the complete kinematic event
  that produced the observed projection. The recognized symbol is
  the mode of the posterior distribution over gestures; the recognized
  meaning is the mode of the posterior over intentions. Classification
  is the special case where the posterior is sharply peaked at a single
  gesture class.
\end{namedthm}

\begin{exercise}
  Write out the full Bayesian inverse problem for reconstructing a
  handwritten letter from a binarized (black/white) image. Define
  $P(O | \gamma)$ and $P(\gamma)$ explicitly for this case.
\end{exercise}

\begin{exercise}
  In the constraint multiplication formula, show that adding an
  independent observation can never \emph{increase} the entropy
  of the posterior (i.e., never makes reconstruction less precise).
  Under what conditions does adding an observation leave the entropy
  unchanged?
\end{exercise}

\begin{exercise}[title={Harder}]
  The latent intention model introduces three levels: observation,
  gesture, and intention. Extend this to a four-level model with
  an additional ``goal'' layer above intention (what the person
  is trying to achieve across multiple words or sentences).
  What additional data would be needed to reconstruct goals?
\end{exercise}
